\chapter{Literature Review}
\label{chap:literature_review}

This chapter builds the theoretical and empirical foundation for the study by synthesizing research on financial liberalization, digital transformation, and computational measurement. It examines the main mechanisms through which access to international capital may impact firm behavior, assesses existing evidence on contextual differences, and highlights unresolved gaps in identification and mechanism integration.

\section{Capital Market Liberalization: Theory and Evidence}

\subsection{Theoretical Foundations of Financial Liberalization}

Capital market liberalization generally refers to removing barriers that separate domestic financial systems from international portfolio flows~\cite{Long2024}. In essence, liberalizing equity markets can improve allocative efficiency by enabling capital transfer to higher-productivity companies, lowering financing costs through more investor participation, and increasing risk-sharing via diversified ownership structures~\cite{Zhu2023}. The traditional welfare argument suggests that firms facing financial constraints should benefit more when foreign capital is available~\cite{Yin2023}. In segmented markets, limitations of domestic savings, home bias, and informational frictions can sustain high costs linked to external finance. Liberalization efforts can reduce these frictions by boosting competition among investors and increasing market liquidity~\cite{HLiu2024}.

Nevertheless, theoretical frameworks also highlight non-linearities. The benefits linked to liberalization depend on the quality of institutions, the effectiveness of legal enforcement, disclosure standards, and the strength of market infrastructure~\cite{Stiglitz2004}. In environments with weak institutional frameworks, liberalization may increase market volatility or shift resources toward short-term opportunities instead of encouraging productive long-term investments~\cite{JXZhang2023}. From a corporate perspective, liberalization can be seen as a change in the optimization problems facing managers and controlling shareholders~\cite{ChengMasron2023}. It affects the relative cost of external financing, the expected investment horizon of investors, and the level of monitoring for capital provision~\cite{Chen2025}.

The MSCI inclusion mechanism is particularly informative because it combines signaling and flow effects. Inclusion can automatically trigger passive fund allocations while also signaling quality and accessibility to active investors.

\subsection{Effects on Corporate Investment, Innovation, and Governance}

Empirical studies generally find that market opening influences corporate behavior through investment and governance channels~\cite{Tian2023}. On the investment side, reduced financing frictions can broaden firms' feasible project options~\cite{RZhang2023}. Firms may increase R\&D, capital expenditure, and strategic intangible investments when expected financing constraints lessen~\cite{Yin2023}. On the governance side, foreign institutional investors can strengthen monitoring and raise transparency expectations.

Evidence from prior work supports this logic. Chen et al.~\cite{Chen2022} show that digital transformation is associated with higher analyst coverage and improved information accuracy. Du and Wang~\cite{DuWang2024} provide direct evidence that financial support strengthens the innovation effects of digital transformation, particularly in environments with tighter conventional financing constraints. Chen and Zhang~\cite{ChenZhang2024} and Feng et al.~\cite{Feng2024} extend this argument to firms' financing structures.

\subsection{Evidence from Emerging Market Contexts}

Emerging markets provide valuable evidence because they combine high growth potential with institutional frictions that highlight financing channels. Studies in these markets often detect stronger real effects of financial reforms compared to more developed markets where constraints are already less severe~\cite{StiglitzOcampo2008}. The Chinese evidence is especially useful due to detailed data and specific policy information~\cite{Yan2023}. However, a direct causal link from liberalization shocks to firm-level digital transformation outcomes is less developed.

Zhang et al.~\cite{Zhang2024} show that contextual heterogeneity is central: estimated effects vary with regional market development, ownership structure, and firm size. This pattern implies that pooled average treatment estimates may mask economically meaningful distributional differences. Stiglitz~\cite{Stiglitz2000} similarly emphasizes that policy shocks interact with organizational absorptive capacity.

\section{Digital Transformation: Conceptual Framework and Determinants}

\subsection{Defining Digital Transformation: Technological and Organizational Dimensions}

Digital transformation extends beyond just acquiring technology. It involves continuous, firm-wide reconfiguration of value creation, operational processes, and organizational coordination through digital technologies~\cite{Yang2023}. A narrow view that only considers software or hardware spending misses the strategic and organizational aspects that influence whether digital investments lead to productivity improvements. Conceptually, digital transformation includes at least three layers~\cite{Verhoef2021,XueZhang2024}: infrastructure adoption, process transformation, and organizational transformation.

These layers depend on each other. Infrastructure without process integration can yield limited results~\cite{WangLi2024}. Process changes without organizational incentives can prevent progress. Organizational commitment without funding can remain superficial~\cite{XueZhang2024}. Therefore, digital transformation is best viewed as a coordinated investment journey rather than a single technological event~\cite{Chanias2019}.

\subsection{Core Technology Domains: AI, Big Data, Cloud Computing, IoT, and Blockchain}

The technology areas highlighted in this thesis align with key trends in modern corporate digitalization. Artificial intelligence and machine learning support prediction, optimization, and automation across both front-end and back-end operations~\cite{Babina2024}. Big data systems enable companies to turn large, fast-moving information into operational insights. Cloud computing offers scalable infrastructure that reduces fixed ICT costs~\cite{Adeola2020}. Internet of Things technologies link physical assets to digital monitoring and control systems~\cite{Kommisetty2022}. Blockchain and distributed ledger technology provide traceability, tamper resistance, and coordination across organizational boundaries~\cite{Ansari2025}. Enterprise digitization initiatives combine these technologies with ERP upgrades, process engineering, and platform strategies.

These domains vary in capital intensity, knowledge requirements, and implementation risks~\cite{Wang2024}. Existing research supports this multidimensional view~\cite{Ren2024}.

\subsection{Drivers, Barriers, and Firm-Level Heterogeneity}

Drivers of digital transformation include financial slack, managerial capability, data infrastructure, and complementary human capital. Barriers are also significant. High uncertainty about benefits can discourage investment, particularly in firms under pressure for short-term performance. Talent shortages in data engineering, AI deployment, and cybersecurity limit progress~\cite{ZhangCao2023}. Legacy systems can increase integration costs.

These barriers imply heterogeneous response capacities. Larger firms may benefit from economies of scale in technology investment~\cite{Gao2024}. Companies in developed regions might access more extensive supplier networks and skilled labor. Non-state firms often face greater market competition, creating stronger incentives to adopt productivity-enhancing digital strategies~\cite{Jin2023}.

\subsection{Digital Transformation in Emerging Markets: Opportunities and Constraints}

Emerging markets offer a unique digitalization landscape with both growth opportunities and structural challenges~\cite{Zhao2023}. In these situations, access to external capital markets can serve a catalytic function. International investors may not only supply funding but also bring governance oversight, benchmarking standards, and a focus on long-term capability building~\cite{Li2025}. Specifically for China, the coexistence of advanced digital ecosystems and uneven regional development creates a particularly revealing context~\cite{Du2024}.

\section{Mechanisms Linking Capital Market Liberalization to Digital Transformation}

\subsection{The Financing Constraint Channel}

The financing-constraint mechanism is the most direct way connecting liberalization to digital transformation. Digital projects often require large upfront expenses in software systems, data architecture, cloud migration, machine learning platforms, and organizational training~\cite{Kose2003}. By broadening access to equity investors and lowering the returns they expect, liberalization can reduce the user cost of capital for such projects~\cite{LiPang2023}.

Existing studies strongly support this channel.~\cite{Hordofa2024} show that capital-market financial support significantly enhances the innovation effects of digital transformation.~\cite{YLiu2024} find lower cost of equity among firms undergoing digital transformation. These results suggest that financing and digitalization can reinforce each other. The channel is unlikely to be consistent across all firms. Firms with initially binding financing constraints should show larger responses~\cite{Huang2023}.

\subsection{Corporate Governance and International Investor Monitoring}

A second mechanism is governance upgrading through international investor monitoring. Foreign institutional investors often demand higher disclosure quality, stronger board accountability, and clearer strategic communication~\cite{Chen2025}. Under these conditions, managers may have stronger incentives to focus on long-term value-enhancing projects and reduce agency-driven underinvestment.

~\cite{Ma2025} provide relevant supporting evidence by showing that digital transformation is linked to improved information environments through increased analyst coverage.~\cite{CLi2023} show that governance effects may also interact with ownership structure.

\subsection{Knowledge Spillovers and Global Best Practice Diffusion}

A third mechanism operates through information and knowledge spillovers. Capital market liberalization boosts interactions between domestic firms and global investors, analysts, and intermediaries~\cite{Zheng2021,Zhu2023}. Through investor meetings, analyst reports, benchmarking exercises, and governance engagement, firms can access implementation templates and performance standards~\cite{Sun2022}. This mechanism aligns with broader themes in the literature~\cite{HLiu2024,LiShen2021}.

\subsection{Market Signaling, Certification, and Reputational Effects}

Liberalization can impact digital transformation through signaling and certification processes. Inclusion in a major global benchmark is often seen as recognition of minimum standards in accessibility, liquidity, and market relevance~\cite{CLi2023}. This signal can lower external uncertainty about a firm's quality.~\cite{Dong2023} discuss distinguishing between substantive and symbolic effects.

\section{Text-Based Measurement of Firm Characteristics}

\subsection{Computational Approaches in Corporate Finance and Economics}

The rise of text-as-data methods has revolutionized empirical corporate research by enabling high-frequency, large-sample analysis~\cite{Yang2024}. In the era of digital transformation, computational text methods are especially valuable because direct accounting measures are often incomplete~\cite{Bellstam2017}. Methodologically, text-based approaches enhance scalability and consistency compared to purely manual coding~\cite{Ahci2023}. However, computational methods need careful validation~\cite{Zimmermann2023}.

\subsection{NLP-Based Measurement of Innovation and Digital Activity}

Afzal~\cite{Afzal2024} shows that Chinese-language annual reports can provide credible NLP-based indicators of digital activity. Kumar~\cite{Kumar2024} further argues for domain-cluster keyword extraction rather than generic buzzword counts. Li et al.~\cite{YLi2023} position NLP-based measurement as a complement to, rather than a substitute for, financial indicators.

\subsection{Validation Strategies for Text-Derived Indices}

Validation is central to any text-based empirical strategy. This thesis applies three validation layers: content validity, construct validity, and robustness validity. Content validity is established through careful dictionary design and expert review. Construct validity is evaluated by comparing text-derived indices with related observable variables. Robustness validity is assessed through various alternative outcome structures.

\subsection{Challenges in Chinese-Language Disclosure Environments}

Chinese-language textual analysis presents technical and interpretive challenges~\cite{YLi2023,Wang2022}. Chinese text segmentation is complex because characters are written without spaces. Another challenge is variation in strategic disclosure~\cite{Yang2024}. A third challenge involves distinguishing between intention and implementation.

Despite these challenges, Chinese annual reports remain some of the most comprehensive and standardized data sources for long-term firm-level analysis. With careful design and validation, they offer a strong empirical foundation for studying digital transformation at scale.

\section{Research Gaps and Positioning of This Study}

The reviewed literature reveals four core gaps that motivate this thesis.

\begin{itemize}
    \item \textbf{Linkage gap}: Research on capital market liberalization and research on digital transformation have mostly developed independently.
    \item \textbf{Identification gap}: Many studies rely on cross-sectional correlations or panel models that may face endogeneity issues.
    \item \textbf{Mechanism integration gap}: Existing evidence suggests financing, governance, and knowledge channels are all plausible but often examined in isolation.
    \item \textbf{Measurement scalability gap}: Direct digital investment metrics are often incomplete in emerging-market accounting systems.
\end{itemize}

This thesis addresses these gaps by examining the MSCI A-share inclusion event as a near-natural experiment to determine the causal effect of liberalization on digital transformation. The literature synthesis provides direct support for this positioning. What is missing is a design that integrates these insights into a single coherent causal framework. This thesis is constructed to provide that framework.
